\section{Prompt Evolution}
\label{sec:prompt-evolution}

This appendix presents the human-authored initial prompts alongside their autonomously optimized counterparts for all five field categories: the core fields (8 fields) and four extended field groups. For the four extended field groups, we embed the representative top-performing prompts (M8, inference=Deepseek-70B, critique=GPT-oss-120B). Across these prompts, the optimizer generally learned to: (i) tightly scope the extraction window to the main article while ignoring cited abstracts/supplementary fluff; (ii) normalize units and apply explicit defaults (thickness/time/concentration) instead of heuristic guessing; (iii) disambiguate primary device elements from auxiliary layers (e.g., gate dielectric vs. encapsulation vs. reference electrode); (iv) collapse multiple material or device variants into a single representative record when the paper only tweaks minor formulations; (v) extract numeric values directly with required conversions (e.g., pH$\rightarrow$[H$^+$], ppm from P/P$_v$, seconds from minutes); and (vi) forbid inference---fields remain "" unless explicitly stated. These behaviors collectively reduce hallucination, enforce consistency, and preserve the most representative device configuration per paper.

Concrete improvements observed across the five prompt families include: (1) \textit{Scope control}---original prompts often swept in cited abstracts or supplementary blurbs; optimized prompts now restrict extraction to the main article body (and, when needed, a specified table or subsection). (2) \textit{Unit rigor}---where initial prompts left units free-form, optimized prompts demand explicit conversions (e.g., minutes$\rightarrow$seconds for response/recovery times; P/P$_v$$\rightarrow$ppm; pH$\rightarrow$[H$^+$] for bounds) or defaults when absent. (3) \textit{Device disentangling}---early instructions blurred substrate/dielectric/encapsulation; optimized prompts explicitly separate gate dielectric from protective coatings and require empty strings when not stated. (4) \textit{Variant consolidation}---initial versions invited multiple records for minor material tweaks; optimized prompts collapse variants into a single representative record unless the paper truly reports distinct devices. (5) \textit{Field-specific cues}---material-layer prompts add dopant handling and thickness defaults; sensitivity/response prompts require table parsing and unit rounding; synthesis/thermal prompts stress atmosphere/time/temperature capture while ignoring non-functional processing layers. Collectively, these edits yield cleaner, reproducible JSON with fewer missing-unit errors and no fabricated values.

Table~\ref{tab:extraction_fields} summarizes the 28 extraction fields organized into five prompt categories.

\begin{table*}[htbp]
\centering
\caption{Summary of LLM extraction fields by prompt category. A total of 28 fields are extracted across five prompt families.}
\label{tab:extraction_fields}
\small
\begin{tabular}{lcp{12cm}}
\toprule
\textbf{Category} & \textbf{Fields} & \textbf{Field Names} \\
\midrule
Core & 8 & sensor\_type, detect\_target, lower\_detection\_limit, upper\_detection\_limit, probe\_material, test\_operating\_temperature, pH\_value, test\_medium \\
\addlinespace
Electrode Architecture & 4 & gate, source, drain, structure\_design\_type \\
\addlinespace
Material Layer & 7 & substrate, substrate\_thickness, channel, dielectric\_layer, dielectric\_layer\_thickness, surface\_functionalization, structure\_dimensionality \\
\addlinespace
Sensitivity/Response & 4 & response\_time, recovery\_time, sensitivity\_numerator, sensitivity\_denominator \\
\addlinespace
Synthesis/Thermal & 5 & annealing\_temperature, annealing\_time, annealing\_atmosphere, hydrothermal\_temperature, hydrothermal\_time \\
\midrule
\textbf{Total} & \textbf{28} & \\
\bottomrule
\end{tabular}
\end{table*}
